% Practice Test 21 — Virginia Edition
% Questions: 30 (24 MC + 6 SA)
% Generated by scripts/generate_practice_tests.py

\begin{practiceQuestion}{1-4-q19}{sa}
\begin{questionText}
Write an equation for the proportional relationship where $3$ tickets cost $\$22.50$.
\end{questionText}
\correctAnswer{$c = 7.50t$}
\explanation{$k = \frac{22.50}{3} = 7.50$ dollars per ticket. The equation is $c = 7.50t$.}
\end{practiceQuestion}

\begin{practiceQuestion}{2-1-q05}{mc}
\begin{questionText}
What is $40\%$ of $150$?
\end{questionText}
\choiceA{$45$}
\choiceB{$50$}
\choiceC{$55$}
\choiceD{$60$}
\correctAnswer{D}
\explanation{$40\% = 0.40$. Multiply: $0.40 \times 150 = 60$.}
\end{practiceQuestion}

\begin{practiceQuestion}{2-2-q15}{mc}
\begin{questionText}
Which equation is equivalent to the proportion $\dfrac{n}{250} = \dfrac{16}{100}$?
\end{questionText}
\choiceA{$n = 250 \div 16$}
\choiceB{$n = 16 \times 250$}
\choiceC{$100n = 250 \times 16$}
\choiceD{$16n = 250 \times 100$}
\correctAnswer{C}
\explanation{Cross-multiply: $100 \times n = 250 \times 16$, which gives $100n = 4000$.}
\end{practiceQuestion}

\begin{practiceQuestion}{2-3-q03}{mc}
\begin{questionText}
A bike costs $\$200$. It is on sale for $35\%$ off. What is the sale price?
\end{questionText}
\choiceA{$\$70$}
\choiceB{$\$130$}
\choiceC{$\$135$}
\choiceD{$\$165$}
\correctAnswer{B}
\explanation{Multiply by the decrease multiplier: $200 \times 0.65 = \$130$.}
\end{practiceQuestion}

\begin{practiceQuestion}{2-7-q12}{mc}
\begin{questionText}
Two students estimated the number of beans in a jar. The actual count was $200$. Student A guessed $180$; Student B guessed $220$. Who had the greater percent error?
\end{questionText}
\choiceA{Student A}
\choiceB{Student B}
\choiceC{Same percent error}
\choiceD{Cannot be determined}
\correctAnswer{C}
\explanation{A: $\frac{|180-200|}{200} = 10\%$. B: $\frac{|220-200|}{200} = 10\%$. Both have $10\%$ error.}
\end{practiceQuestion}

\begin{practiceQuestion}{3-2-q04}{mc}
\begin{questionText}
What is $9 + (-14)$?
\end{questionText}
\choiceA{$23$}
\choiceB{$-23$}
\choiceC{$-5$}
\choiceD{$5$}
\correctAnswer{C}
\explanation{Different signs: $14 - 9 = 5$. The larger absolute value ($14$) is negative, so the answer is $-5$.}
\end{practiceQuestion}

\begin{practiceQuestion}{3-3-q08}{mc}
\begin{questionText}
What is the distance between $-9$ and $4$ on a number line?
\end{questionText}
\choiceA{$5$ units}
\choiceB{$-5$ units}
\choiceC{$13$ units}
\choiceD{$-13$ units}
\correctAnswer{C}
\explanation{Distance $= |{-9} - 4| = |{-13}| = 13$ units. Distance is always positive.}
\end{practiceQuestion}

\begin{practiceQuestion}{3-4-q14}{mc}
\begin{questionText}
What is $\dfrac{3}{4} - 1\dfrac{1}{2}$?
\end{questionText}
\choiceA{$\dfrac{3}{4}$}
\choiceB{$-\dfrac{3}{4}$}
\choiceC{$\dfrac{1}{4}$}
\choiceD{$-\dfrac{1}{4}$}
\correctAnswer{B}
\explanation{$\frac{3}{4} - \frac{3}{2} = \frac{3}{4} - \frac{6}{4} = -\frac{3}{4}$.}
\end{practiceQuestion}

\begin{practiceQuestion}{4-3-q03}{mc}
\begin{questionText}
Expand $-2(3x + 7)$.
\end{questionText}
\choiceA{$-6x + 7$}
\choiceB{$-6x - 14$}
\choiceC{$-6x + 14$}
\choiceD{$6x - 14$}
\correctAnswer{B}
\explanation{$-2 \times 3x = -6x$ and $-2 \times 7 = -14$. Result: $-6x - 14$.}
\end{practiceQuestion}

\begin{practiceQuestion}{4-4-q06}{mc}
\begin{questionText}
Factor $-9x - 6$.
\end{questionText}
\choiceA{$-3(3x - 2)$}
\choiceB{$-3(3x + 2)$}
\choiceC{$3(-3x - 2)$}
\choiceD{$-9(x + 6)$}
\correctAnswer{B}
\explanation{GCF is $3$. Factor out $-3$: $-9x \div (-3) = 3x$ and $-6 \div (-3) = 2$. Result: $-3(3x + 2)$. Check: $-3(3x + 2) = -9x - 6$ \checkmark.}
\end{practiceQuestion}

\begin{practiceQuestion}{5-4-q18}{mc}
\begin{questionText}
Solve $-0.2n + 5.6 = 4$.
\end{questionText}
\choiceA{$n = 8$}
\choiceB{$n = -8$}
\choiceC{$n = 48$}
\choiceD{$n = -48$}
\correctAnswer{A}
\explanation{Subtract $5.6$: $-0.2n = -1.6$. Divide by $-0.2$: $n = 8$. Check: $-0.2(8) + 5.6 = -1.6 + 5.6 = 4$ \checkmark.}
\end{practiceQuestion}

\begin{practiceQuestion}{5-5-q26}{gmc}
\begin{questionText}
Which inequality is represented by the number line below?

\begin{center}
\begin{tikzpicture}[scale=0.7]
  \draw[thick,->] (-1,0) -- (10,0);
  \foreach \x in {0,...,9} \draw (\x,0.15) -- (\x,-0.15) node[below] {$\x$};
  \draw[thick,funBlue,->] (5,0) -- (10,0);
  \draw[thick,funBlue,fill=white] (5,0) circle (4pt);
\end{tikzpicture}
\end{center}
\end{questionText}
\choiceA{$x \geq 5$}
\choiceB{$x > 5$}
\choiceC{$x \leq 5$}
\choiceD{$x < 5$}
\correctAnswer{B}
\explanation{An open circle at $5$ with shading to the right means $x > 5$ (not including $5$).}
\end{practiceQuestion}

\begin{practiceQuestion}{5-6-q25}{sa}
\begin{questionText}
Solve $7x - 3 \geq 18$. Give the smallest integer that is a solution.
\end{questionText}
\correctAnswer{$x = 3$}
\explanation{Add $3$: $7x \geq 21$. Divide by $7$: $x \geq 3$. The smallest integer solution is $3$.}
\end{practiceQuestion}

\begin{practiceQuestion}{6-1-q09}{mc}
\begin{questionText}
A model car uses a scale of $1:24$. The real car is $180$ inches long. How long is the model?
\end{questionText}
\choiceA{$7.5$ inches}
\choiceB{$15$ inches}
\choiceC{$24$ inches}
\choiceD{$156$ inches}
\correctAnswer{A}
\explanation{Divide the actual length by the scale factor: $180 \div 24 = 7.5$ inches.}
\end{practiceQuestion}

\begin{practiceQuestion}{6-2-q10}{mc}
\begin{questionText}
A blueprint uses $1$ in $= 12$ ft. You redraw at $1$ in $= 4$ ft. A door that is $0.5$ in wide on the original becomes how wide?
\end{questionText}
\choiceA{$0.5$ in}
\choiceB{$1$ in}
\choiceC{$1.5$ in}
\choiceD{$3$ in}
\correctAnswer{C}
\explanation{Scale ratio: $\frac{12}{4} = 3$. New width: $0.5 \times 3 = 1.5$ in.}
\end{practiceQuestion}

\begin{practiceQuestion}{6-5-q13}{mc}
\begin{questionText}
A hexagonal prism is sliced parallel to its base. What is the cross-section?
\end{questionText}
\choiceA{Rectangle}
\choiceB{Triangle}
\choiceC{Hexagon}
\choiceD{Circle}
\correctAnswer{C}
\explanation{Any cut parallel to the base of a prism gives a shape congruent to the base. A hexagonal prism has a hexagonal base, so the cross-section is a hexagon.}
\end{practiceQuestion}

\begin{practiceQuestion}{6-6-q20}{sa}
\begin{questionText}
Find the supplement of $43°$.
\end{questionText}
\correctAnswer{$137°$}
\explanation{$180° - 43° = 137°$.}
\end{practiceQuestion}

\begin{practiceQuestion}{6-7-q24}{sa}
\begin{questionText}
Point $S(-3, -5)$ is reflected over the $x$-axis, then translated right $4$. What are the final coordinates?
\end{questionText}
\correctAnswer{$(1, 5)$}
\explanation{Reflect over $x$-axis: $(-3, -5) \to (-3, 5)$. Translate right $4$: $(-3 + 4, 5) = (1, 5)$.}
\end{practiceQuestion}

\begin{practiceQuestion}{6-8-q06}{mc}
\begin{questionText}
Two similar triangles have a scale factor of $\frac{2}{5}$. If a side of the larger triangle is $30$ cm, what is the corresponding side of the smaller triangle?
\end{questionText}
\choiceA{$6$ cm}
\choiceB{$12$ cm}
\choiceC{$15$ cm}
\choiceD{$75$ cm}
\correctAnswer{B}
\explanation{Smaller side $= 30 \times \frac{2}{5} = 12$ cm.}
\end{practiceQuestion}

\begin{practiceQuestion}{7-4-q25}{sa}
\begin{questionText}
A trapezoidal garden has parallel sides of $8$ m and $12$ m, and a height of $5$ m. What is its area?
\end{questionText}
\correctAnswer{$50$ m$^2$}
\explanation{$A = \frac{1}{2}(b_1 + b_2) \times h = \frac{1}{2}(8 + 12) \times 5 = \frac{1}{2} \times 20 \times 5 = 50$ m$^2$.}
\end{practiceQuestion}

\begin{practiceQuestion}{7-5-q14}{mc}
\begin{questionText}
A rectangular prism is $6$ ft by $6$ ft by $10$ ft. What is its surface area?
\end{questionText}
\choiceA{$72$ ft$^2$}
\choiceB{$240$ ft$^2$}
\choiceC{$312$ ft$^2$}
\choiceD{$360$ ft$^2$}
\correctAnswer{C}
\explanation{$SA = 2(6 \times 6) + 2(6 \times 10) + 2(6 \times 10) = 72 + 120 + 120 = 312$ ft$^2$.}
\end{practiceQuestion}

\begin{practiceQuestion}{7-6-q21}{sa}
\begin{questionText}
A triangular prism has a triangular base with base $10$ m and height $6$ m. The prism is $8$ m long. What is the volume?
\end{questionText}
\correctAnswer{$240$ m$^3$}
\explanation{$B = \frac{1}{2} \times 10 \times 6 = 30$ m$^2$. $V = 30 \times 8 = 240$ m$^3$.}
\end{practiceQuestion}

\begin{practiceQuestion}{7-7-q14}{mc}
\begin{questionText}
A cylindrical container has a volume of $628$ cm$^3$ and a height of $8$ cm. What is the radius? Use $\pi \approx 3.14$.
\end{questionText}
\choiceA{$3$ cm}
\choiceB{$4$ cm}
\choiceC{$5$ cm}
\choiceD{$6$ cm}
\correctAnswer{C}
\explanation{$628 = 3.14 \times r^2 \times 8 = 25.12 \, r^2$. $r^2 = 628 \div 25.12 = 25$. $r = 5$ cm.}
\end{practiceQuestion}

\begin{practiceQuestion}{8-1-q05}{mc}
\begin{questionText}
A random sample means that:
\end{questionText}
\choiceA{Only the best members are chosen}
\choiceB{The largest possible group is chosen}
\choiceC{Every member of the population has an equal chance of being selected}
\choiceD{Only volunteers are selected}
\correctAnswer{C}
\explanation{In a random sample, every member of the population has an equal chance of being chosen.}
\end{practiceQuestion}

\begin{practiceQuestion}{8-3-q12}{mc}
\begin{questionText}
When is a dot plot a better choice than a box plot for comparing two groups?
\end{questionText}
\choiceA{When data sets are very large}
\choiceB{When you want to see every individual data point}
\choiceC{When you only want to compare medians}
\choiceD{When the data is categorical}
\correctAnswer{B}
\explanation{Dot plots show every data point, which is useful for small data sets where individual values matter.}
\end{practiceQuestion}

\begin{practiceQuestion}{8-4-q11}{mc}
\begin{questionText}
Data set: $12, 15, 18, 20, 22, 25, 28$. What is the median?
\end{questionText}
\choiceA{$18$}
\choiceB{$20$}
\choiceC{$22$}
\choiceD{$19$}
\correctAnswer{B}
\explanation{There are $7$ values. The middle value (4th) is $20$.}
\end{practiceQuestion}

\begin{practiceQuestion}{8-5-q14}{mc}
\begin{questionText}
A back-to-back stem-and-leaf plot is used to:
\end{questionText}
\choiceA{Show one data set with two stems}
\choiceB{Compare two data sets using the same stems}
\choiceC{Display data that has been doubled}
\choiceD{Show data in a circle}
\correctAnswer{B}
\explanation{A back-to-back plot has leaves for one group on the left and for another group on the right, sharing the same stems.}
\end{practiceQuestion}

\begin{practiceQuestion}{9-3-q06}{mc}
\begin{questionText}
A bag contains marbles of unknown colors. After $40$ draws (with replacement), red is drawn $14$ times. What is the experimental probability of drawing red?
\end{questionText}
\choiceA{$\frac{14}{40}$}
\choiceB{$\frac{14}{26}$}
\choiceC{$\frac{26}{40}$}
\choiceD{$\frac{40}{14}$}
\correctAnswer{A}
\explanation{$P = \frac{\text{red draws}}{\text{total draws}} = \frac{14}{40} = 0.35$.}
\end{practiceQuestion}

\begin{practiceQuestion}{9-6-q03}{mc}
\begin{questionText}
A coin is flipped and a die is rolled. What is the probability of getting tails and an even number?
\end{questionText}
\choiceA{$\frac{1}{12}$}
\choiceB{$\frac{1}{6}$}
\choiceC{$\frac{1}{4}$}
\choiceD{$\frac{1}{2}$}
\correctAnswer{C}
\explanation{$P(\text{tails}) = \frac{1}{2}$ and $P(\text{even}) = \frac{3}{6} = \frac{1}{2}$. $P = \frac{1}{2} \times \frac{1}{2} = \frac{1}{4}$. Or: $3$ favorable out of $12$ total outcomes.}
\end{practiceQuestion}

\begin{practiceQuestion}{9-7-q01}{mc}
\begin{questionText}
What is a simulation?
\end{questionText}
\choiceA{A way to calculate exact probabilities}
\choiceB{A model that uses random numbers or objects to imitate a real event}
\choiceC{A method that always gives the same result}
\choiceD{A graph that shows all possible outcomes}
\correctAnswer{B}
\explanation{A simulation uses random numbers, coins, dice, or spinners to model a real-world event and estimate probabilities.}
\end{practiceQuestion}
